\chapter{Scientific Interpretation and Research Decisions}
\label{ch:interpretation}

Chapter~\ref{ch:activities} described what was built; Chapter~\ref{ch:results}
established what the integrated evidence shows. This chapter asks what
those findings mean scientifically, and why they justified the major
research decisions the project actually took. It does not repeat the
chronological narrative or the metric-by-metric synthesis of the two
preceding chapters; results are cited by reference, briefly, only where a
specific number is the direct subject of an interpretive claim. This
chapter distinguishes throughout between what the evidence supports
observing and what may be inferred from it, and prefers hedged language
("is consistent with," "suggests," "cannot distinguish between") to
causal claims wherever a causal mechanism is not established by the
project's or the predecessor's own documentation.

\section{From Visible Corrosion to Structural Condition: An Evidence Hierarchy}
\label{sec:interp-hierarchy}

The project's successive research questions form an evidence hierarchy,
not a sequence of independent tasks — and not a literal universal
computational pipeline in which each stage is generated exclusively from
the stage immediately above it. Figure~\ref{fig:evidence-hierarchy} sets
this out as a ladder of increasing inferential distance from directly
observed quantities, decreasing supervision density, and increasing
dependence on modelling assumptions; actual feature inputs at each stage
are phase-dependent, not fixed by the ladder position alone. Direct
observations and labels — images, visible-corrosion labels, and
specimen/exposure metadata — sit at the top, densely available for nearly
every observation (Section~\ref{sec:data-modalities}). Image
representations and image-based estimates are computed from these
observations. Hidden structural condition is a model's estimate, fit on 48
terminal measurements and applied to weeks at which no structural
measurement exists (Section~\ref{sec:structural-supervision}); its actual
inputs vary by phase rather than depending on the image-
representation tier alone — the interpretable-feature phase's structural
models are mixed-input (metadata, manually assigned corrosion labels, and
image descriptors together), while the robustness-quantification phase's
final selections use metadata only (Section~\ref{sec:results-structural}).
Degradation trajectories are curves fitted to those estimated, not
measured, series (Section~\ref{sec:phase3-degradation}). Proxy-RUL is a
threshold-crossing time read off an extrapolation of one of those fitted
curves (Section~\ref{sec:proxy-rul}).

\begin{figure}[htbp]
\centering
\resizebox{0.85\textwidth}{!}{% Report schematic (conceptual diagram, not a measured result / not to
% scale, and not a computational data-flow diagram). Position in this
% ladder reflects increasing inferential distance from directly observed
% quantities, decreasing supervision density, and increasing dependence on
% modelling assumptions -- it does NOT mean that each stage is computed
% exclusively from the stage immediately above it. Actual feature inputs
% are phase-dependent: e.g. hidden structural estimates in one phase are
% mixed-input (metadata + manually assigned corrosion labels + image
% descriptors), while another phase's final structural selections are
% metadata-only. Plain connecting lines (no arrowheads) are used
% deliberately, to avoid implying a universal computational dependency.
\begin{tikzpicture}[
    node distance=9mm and 10mm,
    stage/.style={draw, rounded corners, align=center, minimum width=6.6cm,
                minimum height=1.15cm, font=\small},
    observed/.style={stage, fill=gray!6},
    derived/.style={stage, fill=gray!18},
    screening/.style={stage, fill=gray!30},
    conn/.style={thin, dashed},
    bracket/.style={very thick, -{Latex[length=2.5mm]}},
    sidelbl/.style={font=\scriptsize\itshape, align=left}
]

\node[observed] (obs) {Direct observations / labels\\(images, visible-corrosion labels, metadata)};
\node[observed, below=of obs] (images) {Image representations and image-based estimates\\(handcrafted or learned)};
\node[derived, below=of images] (structural) {Hidden structural estimates\\(model-derived; inputs are phase-dependent)};
\node[derived, below=of structural] (trajectory) {Degradation trajectories\\(curves fitted to model-estimated series)};
\node[screening, below=of trajectory] (rul) {Proxy-RUL / threshold status\\(screening quantity)};

\draw[conn] (obs) -- (images);
\draw[conn] (images) -- (structural);
\draw[conn] (structural) -- (trajectory);
\draw[conn] (trajectory) -- (rul);

\draw[bracket] ([xshift=-9mm]obs.north west) -- ([xshift=-9mm]rul.south west);

\node[sidelbl, left=13mm of structural, text width=2.6cm, align=center]
  {increasing\\inferential\\distance};

\node[sidelbl, right=8mm of images, yshift=-2mm, text width=4.8cm]
  {\textbf{Supervision:} dense $\rightarrow$ sparse\\(every image labelled $\rightarrow$ 48 terminal measurements only)};
\node[sidelbl, right=8mm of trajectory, yshift=-2mm, text width=4.8cm]
  {\textbf{Modelling assumptions:} few $\rightarrow$ many\\(each lower tier adds assumptions not required by the tier above)};

\node[font=\tiny\itshape, align=left, below=6mm of rul, text width=11.5cm]
  {Dashed connectors indicate ladder position (inferential distance), not a\\
   universal computational dependency: actual feature inputs are\\
   phase-dependent (see text).};

\end{tikzpicture}
}
\caption{Conceptual report schematic (not a measured result). Position in
the ladder reflects increasing inferential distance from directly observed
quantities, decreasing supervision density, and increasing dependence on
modelling assumptions — not a universal computational dependency in which
each stage is generated solely from the stage immediately above it; actual
feature inputs are phase-dependent, as discussed in the surrounding text.
This illustrates why evidential confidence decreases downstream, not that
downstream estimation is impossible.}
\label{fig:evidence-hierarchy}
\end{figure}

This hierarchy is offered as an organising explanation for a pattern
already established in Chapter~\ref{ch:results}: results become
progressively less robust and more heavily qualified moving down it.
That pattern is consistent with supervision becoming sparser and model
dependence compounding at each step; it is not, on its own, proof that any
particular downstream quantity cannot in principle be estimated more
reliably with different data or methods. The distinction matters for how
the rest of this chapter should be read: each section below interprets why
the evidence looks the way it does, without treating the hierarchy itself
as an argument for impossibility.

\section{Why Visible Corrosion Was the Strongest Image-Based Task}
\label{sec:interp-visible}

Three independent phases, using unrelated representations (a fixed colour
threshold, a frozen deep embedding, and an explicitly engineered
multi-family feature set), each reached useful grouped-holdout accuracy
for a visible-corrosion target (Section~\ref{sec:results-visible}). This
consistency across unrelated representations is consistent with the
underlying task being genuinely well suited to image-based estimation:
corrosion is a surface phenomenon that manifests directly in the
photographic record, and it is the one target family labelled at the same
fine time resolution as the images themselves
(Section~\ref{sec:data-modalities}), giving every phase far more
supervision to learn from than any structural target could offer.

This should not, however, be read as a uniform success story. Peak rust
provides more meaningful independent evidence of genuine image-driven
estimation than the total-rust benchmarks do, precisely because two
separate phases independently found total/current rust predictions
dominated by a feature closely aligned with, or nearly identical to, the
label itself (Section~\ref{sec:results-visible}). The lesson this
supports is specific and worth stating plainly: high predictive accuracy
is not sufficient evidence of a learned visual concept if the dominant
feature substantially reconstructs the target's own definition. This
chapter treats that lesson as a general one, applicable beyond the two
cases in which it was directly observed, not as a defect specific to
those two phases. Robustness, similarly, does not transfer uniformly:
campaign-level robustness for visible corrosion is a replicated finding,
while treatment-level robustness is not
(Section~\ref{sec:results-robustness}). The difference may partly reflect
changes in treatment-group definition, size, or composition between the
two implementations; the available evidence establishes the disagreement
itself, not its mechanism, and does not rule out the visual task as a
contributor.

\section{Why Structural Inference Was Harder}
\label{sec:interp-structural}

Hidden structural condition differs from visible corrosion along every
axis identified in Chapter~\ref{ch:dataset}: it is measured once per
specimen rather than repeatedly, at a single terminal timepoint rather
than throughout the observation period, for only 48 specimens rather than
791 observations (Section~\ref{sec:structural-supervision}). This sparse,
single-timepoint supervision is a major limitation consistent with the
observed instability of structural models relative to visible-corrosion
models, independent of any claim about the images themselves. It is
compounded by the experimental design's near-perfect alignment of
campaign identity with steel-mesh configuration, chloride concentration,
and ageing duration (Section~\ref{sec:confounding}): a model trained on
all 48 rows has, in effect, only two campaign-level combinations of the
principal design variables (mesh configuration, chloride level, and
terminal duration) to learn from — treatment variation still exists within
each campaign — and evaluating it under leave-one-campaign-out asks it to
extrapolate to a combination of those variables it has never encountered
(Section~\ref{sec:loco}).

The available evidence supports a narrow, specific conclusion, and it is
important not to overstate it in either direction. It would be wrong to
conclude that surface corrosion does not affect structural capacity — no
result here tests that physical claim directly, and Chapter~\ref{ch:results}
explicitly avoided asserting it. The evidence instead supports the
narrower claim that this dataset and this experimental design do not
isolate a robust, incremental, image-based relationship to hidden
structural capacity. The distinction is between a physical relationship,
which the available evidence can neither confirm nor rule out, and
predictive identifiability from this specific dataset, which is what was
actually tested and what the project's results speak to. Reading the
mixed-input character of the interpretable-feature phase's structural
models alongside the metadata-only character of the robustness-
quantification phase's final selections (Section~\ref{sec:results-structural})
reinforces this reading: even when image descriptors are made available to
a structural model, the model does not consistently rely on them once
selection is made robustness-aware.

\section{Metadata Dominance: Signal or Shortcut?}
\label{sec:interp-metadata}

Metadata's dominance of the terminal-load prediction task
(Section~\ref{sec:results-incremental}) admits two readings that hold
simultaneously rather than competing. Steel-mesh configuration, chloride
exposure, ageing duration, and treatment protocol are genuine design and
exposure variables with a plausible physical relationship to structural
response; there is nothing suspect, in itself, about a model finding them
predictive. At the same time, because these same variables are almost
perfectly aligned with campaign identity in this sample
(Section~\ref{sec:confounding}), a pooled model's reliance on them cannot
be distinguished, from the pooled evidence alone, from the model simply
having learned which campaign a specimen came from. Both readings are
consistent with the same pooled result, and the available evidence does
not adjudicate between them from the pooled analysis alone.

This is why metadata should not be described as invalid or as a nuisance
variable to be engineered away: it may well carry real information about
structural response, and a metadata-only baseline is exactly the right
comparison against which to judge whether images add anything further
(Section~\ref{sec:multimodal}). What the mesh-stratified and campaign-
holdout evidence adds (Section~\ref{sec:results-structural}) is a second,
independent check on the same question, without relying on the pooled
correlation at all. The failure to retain predictive quality within mesh
groups and under campaign holdout weakens the evidence that the pooled
metadata relationship is transferable beyond campaign-linked structure; it
does not disprove a genuine physical metadata effect, only the claim that
the pooled result demonstrates one. This is why pooled predictive
performance can be genuinely useful — it is still the strongest available
baseline for this task — while simultaneously being insufficient evidence,
on its own, of a transferable structural relationship.

\section{Why the Research Question Was Refocused}
\label{sec:interp-refocus}

The shift from hidden-damage estimation toward the terminal ultimate-load
refocus (Section~\ref{sec:phase4b}) is best read as a scientifically
motivated narrowing of claim scope, not as a retreat from a failed
experiment. The hidden-damage chain — estimate structural condition at
every observation week, fit a degradation curve to those estimates,
extrapolate to a threshold crossing — compounds a separate layer of
modelling uncertainty at each step (Section~\ref{sec:results-degradation}).
Terminal ultimate load, by contrast, is directly measured for every one of
the 48 specimens at a single, well-defined timepoint
(Section~\ref{sec:experimental-context}), which removes the need for any
intervening model-estimated quantity between the data and the target.

Reframing the question around that directly measured endpoint also made
the resulting experiment more bounded and directly testable. "Do images
add information beyond known design and exposure metadata for a directly
measured terminal endpoint?" is a bounded, answerable question with a
specific comparison built in from the start (Section~\ref{sec:multimodal}).
"Can a sparse hidden state be estimated at every intermediate week and
then extrapolated into a remaining-life estimate?" is not incoherent, but
is substantially harder to validate and diagnose, because disagreement at
the final stage cannot be cleanly attributed to one component without
independent intermediate ground truth. The refocused study's own
design — matched feature-set comparisons, a mesh-stratified check, and a
campaign-holdout stress test, all reported together rather than as a
single isolated accuracy number (Section~\ref{sec:phase4b}) — is
consistent with this being a stronger experimental design, not merely a
narrower one.

\section{Scientific Value of the Negative Results}
\label{sec:interp-negative}

Several of the project's most useful findings are negative, and they
should be read as evidence, not as failures to be minimised. That images
do not provide a reliable pooled improvement over metadata for terminal
ultimate load (Section~\ref{sec:results-incremental}) provides evidence
against an otherwise plausible claim in this dataset, under the tested
pooled setting, rather than merely failing to confirm it. That
every feature-set configuration achieves negative $R^2$ within a single
mesh family (Section~\ref{sec:phase4b}) reveals a gap between absolute
error and explained variance that a less careful evaluation would have
missed entirely. That structural performance collapses under
leave-one-campaign-out in every phase that tested it
(Section~\ref{sec:results-structural}) identifies exactly which evaluation
regime future work in this dataset must budget for. That the refined
proxy-RUL configuration yields a zero future-crossings result — no
projected future threshold crossing among any of the tested cases, even
though many are already crossed at baseline or during observation
(Section~\ref{sec:results-degradation}) — is a more informative outcome
than a handful of plausible-looking but unvalidated future estimates would
have been. That the residual-correction attempt on
top of the metadata-only baseline measurably worsened performance
(Section~\ref{sec:phase4b}) was recorded rather than discarded, which is
itself part of what makes the surrounding positive claims credible. Each
of these findings narrows the space of defensible claims, exposes a
specific confound or data requirement, and motivates a specific next
check — which is what a negative result is for.

\section{What the Proxy-RUL Work Contributes}
\label{sec:interp-rul}

The proxy-RUL pipeline's contribution is methodological and diagnostic,
not prognostic, and this chapter states that distinction directly rather
than qualifying it away. Executing the four-step derivation end to end,
twice, under two independently configured settings
(Section~\ref{sec:results-degradation}), forced an explicit separation
between observed and model-estimated condition series
(Table~\ref{tab:quantity-availability}) that a less disciplined pipeline
could have left implicit. The refined generation's threshold-status
taxonomy — crossed at baseline, crossed during observation, or not crossed
within the horizon — is a more transparent way of reporting the same
underlying uncertainty than a single ambiguous "weeks remaining" figure
would be, and its own zero future-crossings result is itself a diagnostic
finding: the absence of any projected future crossing, together with the
large number of already-baseline-crossed cases, is consistent with
threshold status being strongly influenced by already-elevated baseline
structural estimates, although that decomposition was not itself
quantified. The physical plausibility of those baseline structural
estimates remains, by the project's own account, an unresolved validation
check rather than a settled fact. The project's own supporting
documentation separately and explicitly recommends exporting genuinely
held-out (out-of-fold) structural predictions in place of the current
in-sample refits for any future degradation-modelling work — naming the
relevant risk directly: that degradation and proxy outputs built on
in-sample structural refits can appear more confident than the underlying
structural model's own robustness evidence would justify. This is a
downstream-optimism risk the project's own documentation makes visible,
not one this chapter is introducing for the first time.

None of this amounts to validated remaining-life estimation, and the
project's own dataset does not support it: true RUL validation would
require either a recorded failure-event time or repeated, rather than
single-terminal, structural measurements per specimen
(Section~\ref{sec:proxy-rul}), and neither is present. The pipeline's
value lies in having built, tested, and critically examined the
derivation chain that such data would require — clarifying exactly what
is still missing — rather than in having produced a usable service-life
number.

\section{Interpretation of the Four-Class Classification Direction}
\label{sec:interp-classification}

The current four-class classification effort is a reasonable next
direction on methodological grounds, independent of what a future trained
classifier will show. It targets a directly observable quantity — image to
severity category — rather than a model-derived or extrapolated one,
using an explicit, current schema kept distinct from the historical
five-level scheme used in the earlier dataset build
(Section~\ref{sec:dataset-audit}). Its
class imbalance was characterised before any training was attempted
(Section~\ref{sec:phase5}), its specimen-level grouped split and leakage-
safety runtime check were implemented and verified rather than merely
documented, and its augmentation recipes were chosen or rejected on stated
physical-plausibility grounds rather than convenience. Read together,
these choices show the branch incorporating, from the outset, the
methodological lessons the earlier regression work had to establish
one at a time over several phases (Section~\ref{sec:interp-lessons}) —
which is a reasonable basis for describing it as well prepared, whatever
its eventual accuracy turns out to be.

This should not be read as more than it is. No classifier has been
trained, compared, or evaluated on this data, and no ResNet-50 or Vision
Transformer experiment exists yet in this project
(Section~\ref{sec:results-classification}). The evaluation practice
already specified for whichever training procedure is used next —
macro-averaged F1 and per-class recall in preference to overall accuracy,
given the severity of the class imbalance, and a test partition left
untouched until final evaluation (Section~\ref{sec:metrics}) — is stated
here as the standard this branch should be held to, not as a result
already obtained.

\section{Methodological Lessons}
\label{sec:interp-lessons}

Several lessons recur across the sections above and are worth drawing
together briefly, as a closing statement rather than a scorecard. Evaluation design can matter more than model complexity: the
project's regime comparisons changed which results were credible more
than any change of model family did (Section~\ref{sec:interp-visible},
Section~\ref{sec:eval-regimes}). A grouped holdout is necessary but not
sufficient: it answers a narrower question than transfer to a new
treatment or campaign, and the project's own later phases treated
additional regimes as required, not optional
(Section~\ref{sec:principles}). Accuracy without attention to target
provenance can be actively misleading, as the total-rust and peak-rust
contrast shows directly (Section~\ref{sec:interp-visible}). A metadata-only
baseline is not an optional control but a precondition for any claim that
images add value (Section~\ref{sec:interp-metadata}). Sparse downstream
labels bound what any model, however capable, can be shown to establish,
independent of modelling effort (Section~\ref{sec:interp-structural}).
Negative results are informative and should shape the next question rather
than be minimised or omitted (Section~\ref{sec:interp-negative}). And
where a directly measured endpoint exists, it is preferable to a longer
chain of model-derived proxies for exactly the reason the terminal-load
refocus illustrates: each additional derived step is an additional,
separately uncertain modelling assumption
(Section~\ref{sec:interp-refocus}).

These lessons are offered as this project's own accumulated methodological
position, consistent with the evidence reviewed in
Chapters~\ref{ch:activities} and~\ref{ch:results}, rather than as general
claims about machine learning practice beyond this project's scope.
